\subsection{章末確認問題}
\begin{enumerate}
  \item モデルが「対象の完全な複製」ではなく「目的を持つ抽象表現」である理由を説明せよ。
  \item モデリング原則の三つ、本質をモデル化する、視点を与える、効果的なコミュニケーションを可能にする、を自分の言葉で説明せよ。
  \item 構文、意味論、語用論の違いを、二つの箱を一本の線で結んだ図を例に説明せよ。
  \item 出金処理を例に、事前条件、事後条件、不変条件をそれぞれ一つずつ挙げよ。
  \item 構造モデルと振る舞いモデルの違いを説明せよ。
  \item モデルの完全性、整合性、正しさは、それぞれ何を確認する観点か。
  \item 追跡可能性が、変更時の影響分析に役立つ理由を説明せよ。
  \item 形式手法が有効になりやすい状況と、導入時の注意点を説明せよ。
  \item プロトタイピングが、通常の開発方法と異なり「よく理解されていない部分」から始める理由を説明せよ。
  \item RAD、XP、Scrum、FDD、Leanの違いを、目的または特徴的な実践に注目して説明せよ。
\end{enumerate}

\begin{pointbox}{解答の方向性}
1は「必要な側面だけを残して複雑さを下げる」と説明できる。2は「何を省くか、どの視点で見るか、誰に伝えるか」を軸に考える。3は「文法、意味、文脈」の区別である。4は「実行前に必要な条件、実行後に保証される条件、前後で変わらない条件」である。5以降は、本文の「何があるか」「どう動くか」「抜け漏れ」「矛盾」「要求適合」「成果物間リンク」「動的協調」という語を使うと説明しやすい。
\end{pointbox}
